\section*{本讲先修与使用方法}
本讲把普通函数推广为“一个输入对应一整套允许结果”的\textbf{对应（correspondence）}。例如预算变化时，可买的商品组合组成一个集合。全文用 $\varphi:S\rightrightarrows T$ 表示对应，$\varphi(x)\subseteq T$ 是它在 $x$ 处的值。除特别讨论空集的页外，默认值非空。$x_n\to x$ 表示输入越来越接近 $x$；$y_n\in\varphi(x_n)$ 表示每次都选择一个当时允许的结果。上标 $q$ 在原课件里通常只是序列编号，不是乘方。

\paragraph{四个基础词。}开集的每个点周围都有一个小球仍在集合中；闭集包含其所有收敛序列的极限；在有限维欧氏空间，紧集等价于闭且有界，也等价于每个序列都有收敛子列且极限仍在集合中；邻域是包含该点某个开球的集合，本文可直接取开邻域。$\varphi(A)=\bigcup_{x\in A}\varphi(x)$ 是所有输入在 $A$ 中时能出现的结果。

\paragraph{先分清三种性质。}闭图：输入和结果\emph{已经一起收敛}时，极限仍然可行。上半连续（开集定义）：包含当前全部结果的开集，也会包含附近的全部结果。下半连续：当前的每个结果，附近都能找到接近它的可行结果。闭图与上半连续并非一般等价；第4页的共同紧值域条件使它们等价。课件前后切换定义，读题时必须先写清假设。第84页把一个“必能提取子列”的强性质称为 o.s.c.，本讲按课件使用这一名称并明确提醒它与常见术语的差异。

\paragraph{建议学习顺序。}先学原页4--14的定义与序列证明，再学15--21的最大值定理；23--53补足开集语言及紧性；54--95训练运算与纠错；96--113系统整理最优化。原页每页均完整保留；下方文字是针对该页的中文说明。课件中的错误论证不能作为考试答案，标有“纠正”的地方给出可用版本。参考背景为 Debreu, \emph{Theory of Value} 第1章第8节；后半部分的推广与纠错是学习补充，不冒充教材原话。

\Slide{1}{从连续函数走向连续对应}
标题页确定主题：连续对应。普通函数 $f(x)$ 是一个点，对应 $\varphi(x)$ 是一个集合。若一名消费者有多个并列最优选择，用函数会丢失信息，对应能把这些选择全部保留下来。本讲的主问题是：环境轻微变化时，可行集合、最高收益、最优选择分别怎样变化？

\Slide{2}{本讲在全章中的位置}
集合与映射提供语言，实数与极限提供收敛，连续函数提供“带入极限”的规则；本讲把这些规则用于集合值映射。下一讲的凸性和最后的固定点定理将使用本讲的上半连续对应。现阶段务必熟悉 $\in$ 与 $\subseteq$：$y\in\varphi(x)$ 是点属于集合，而 $\varphi(x)\subseteq G$ 是整套结果都在 $G$ 中。

\Slide{3}{经济学为什么需要对应}
环境 $x$ 可以包含价格、收入和技术；$\varphi(x)$ 是这一环境下允许的行动集合。例如 $\varphi(p,w)=\{y\ge0:p\cdot y\le w\}$ 是预算集。研究连续性不是要求集合内每个点保持不变，而是要求变化遵守某种“可逼近”的规律。非空性、紧性与连续性在优化中承担不同职责：有选择、能取到最优、最优行为对参数稳定。

\Slide{4}{本节共同假设：值域紧}
课件先取 $S\subseteq\R^K$、\textbf{紧集} $T\subseteq\R^L$，且 $\varphi(x)\subseteq T$。这意味着所有可能结果都装在同一个闭有界容器里。若 $y_n\in\varphi(x_n)$，不管 $x_n$ 如何变化，总可从 $y_n\in T$ 抽出收敛子列。这一点比“每个 $\varphi(x)$ 单独紧”强：$\varphi(x)=\{1/x\}$（$x>0$）每个值都紧，但 $x\downarrow0$ 时结果会逃到无穷远。后续不用共同紧 $T$ 时，需要真正的上半连续加紧值性补上控制。

\Slide{5}{连续性拆成两个方向}
上半连续控制“不能在附近出现脱离当前集合的结果”；下半连续控制“当前结果不能在附近突然无法逼近”。两者合起来才是连续对应。不要把“上、下”理解为图像往上或往下，它们约束的是整套集合变化的不同方向。最优解对应往往只有上半连续，最优值却仍可连续，这正是本讲的核心现象。

\Slide{6}{课件的上半连续：闭图式序列条件}
课件条件逐字读作：\emph{对任意} $x_n\to x_0$，\emph{对任意} 满足 $y_n\in\varphi(x_n)$ 且 $y_n\to y_0$ 的选择，都有 $y_0\in\varphi(x_0)$。结果序列收敛是前提，不能自行假定。它就是闭图的序列条件；在第4页共同紧 $T$ 和闭值背景下，等价于后面的开集式上半连续。

距离说法也可验证：若 $d(y_0,\varphi(x_n))\to0$，选 $y_n\in\varphi(x_n)$ 使
\[\norm{y_n-y_0}<d(y_0,\varphi(x_n))+1/n.\]
右边趋于零，故闭图条件给 $y_0\in\varphi(x_0)$。为什么加 $1/n$？距离是下确界，未必已取到最小值，近似选择总能做到。

\Slide{7}{图2：放大与缩小的区别}
图的横轴是输入，纵向截面才是 $\varphi(x)$，阴影表示图形 $G(\varphi)$，不要把它当普通函数曲线。左图在 $x_0$ 的截面 $[a^1,a^2]$ 可以比邻近截面更大：附近的极限仍落在这一大截面内，因而可以上半连续，但大截面中新增的点未必能从附近逼近。右图在 $x_0$ 只保留较小截面 $[b^1,b^2]$：其中每点能从附近找到近似点，可以下半连续，但邻近多出来的点的极限可能不在截面中。端点的粗细和是否包含是判断闭图的关键。

\Slide{8}{下半连续：为指定目标构造逼近}
这里量词是：任给 $x_n\to x_0$，再任给 $y_0\in\varphi(x_0)$，\emph{存在} 一列 $y_n\in\varphi(x_n)$ 使 $y_n\to y_0$。与第6页区别是，现在目标点先给定，我们负责选择合适的 $y_n$；不是说所有选择都趋向它。

例如常值对应 $\varphi(x)=[0,1]$ 下半连续，因为对任意 $y_0\in[0,1]$ 直接取 $y_n=y_0$。下半连续与 $d(y_0,\varphi(x_n))\to0$ 等价：一方向用 $d\le\norm{y_n-y_0}$；另一方向按第6页加 $1/n$ 近似取点。距离趋零不要求 $y_0$ 本身在每个邻近集合内。

\Slide{9}{连续对应的定义}
连续对应等于同时上半连续和下半连续。一个基本例子是 $\varphi(x)=[a(x),b(x)]$，其中 $a,b$ 连续且 $a\le b$。若 $y_n\in[a(x_n),b(x_n)]$ 且 $y_n\to y$，取极限得 $a(x)\le y\le b(x)$；若指定 $y\in[a(x),b(x)]$，取 $y_n=\min\{b(x_n),\max\{a(x_n),y\}\}$，则 $y_n$ 可行且趋于 $y$。两方向分别验证，才算证明完连续。

\Slide{10}{单值对应退化为普通连续函数}
设 $\varphi(x)=\{f(x)\}$。下半连续中的可选点只有 $y_n=f(x_n)$，故下半连续恰好要求 $f(x_n)\to f(x_0)$。对课件的闭图式上半连续，还要用紧 $T$：若 $f(x_n)$ 不趋于 $f(x_0)$，存在 $\eps>0$ 和子列始终距离至少 $\eps$；紧性再抽出 $f(x_{n_k})\to y$，闭图迫使 $y=f(x_0)$，矛盾。

去掉共同紧值域，闭图不够：$f(0)=0$、$f(x)=1/x$（$x\ne0$）图闭，但在0不连续。若直接采用第24页开集定义，单值上半连续与普通连续的等价不需紧性。

\Slide{11}{上半连续与闭图：双向证明}
图形定义为 $G(\varphi)=\{(x,y):x\in S,\ y\in\varphi(x)\}$，闭是相对于 $S\times T$ 而言。若图闭，取 $x_n\to x_0$、$y_n\to y_0$ 且 $y_n\in\varphi(x_n)$；于是 $(x_n,y_n)\in G$ 且收敛至 $(x_0,y_0)$，闭性给 $(x_0,y_0)\in G$，翻译回去就是 $y_0\in\varphi(x_0)$。

反之，取图中任意收敛序列 $(x_n,y_n)\to(x_0,y_0)\in S\times T$。乘积空间收敛等价于每个坐标收敛，课件序列条件给 $y_0\in\varphi(x_0)$，所以极限仍在图中，图闭。\textbf{考试提醒：}若题目中的 u.s.c. 用开集定义，引用此结论时要同时说明紧值域或适当紧性假设。

\Slide{12}{连续函数与对应的复合}
定义 $\psi=\varphi\circ f$，即先由 $x$ 得到点 $f(x)$，再得到集合 $\varphi(f(x))$。上半连续的证明是把中间变量写出：$x_n\to x$ 推出 $z_n=f(x_n)\to f(x)=z$；若 $y_n\in\psi(x_n)=\varphi(z_n)$ 且 $y_n\to y$，由 $\varphi$ 的条件得 $y\in\varphi(z)=\psi(x)$。

下半连续：任取 $y\in\psi(x)=\varphi(z)$；由 $z_n\to z$，可选 $y_n\in\varphi(z_n)$ 且 $y_n\to y$，这些点正属于 $\psi(x_n)$。\textbf{纠正：}课件第二段开头的 $y\in\varphi(x)$ 应为 $y\in\psi(x)$；$x$ 与 $f(x)$ 所在空间不同，不能混用。

\Slide{13}{有限个对应的笛卡尔积}
$\varphi(x)=\prod_{i=1}^N\varphi_i(x)$ 的元素是向量 $(y_1,\ldots,y_N)$，每个坐标独立满足 $y_i\in\varphi_i(x)$。这常用于多个消费者同时选择：每人一个可行集，社会的可行选择组合是它们的积。有限积的连续性可逐坐标证明；第4页框架中，各个 $T_i$ 紧则乘积也紧。有限维与有限个坐标使“各坐标收敛”等价于“整个向量收敛”。

\Slide{14}{积对应的证明逐步写}
上半连续：任取 $x_q\to x$、$y^q\in\varphi(x_q)$、$y^q\to y$。对每个 $i$，有 $y_i^q\in\varphi_i(x_q)$ 且 $y_i^q\to y_i$；上半连续得 $y_i\in\varphi_i(x)$，全部坐标同时成立即 $y\in\prod_i\varphi_i(x)$。

下半连续：给定 $y=(y_1,\ldots,y_N)\in\varphi(x)$。对每个 $i$，下半连续允许选 $y_i^q\in\varphi_i(x_q)$ 趋于 $y_i$；合成 $y^q=(y_1^q,\ldots,y_N^q)$，它可行，且 $\norm{y^q-y}^2=\sum_i|y_i^q-y_i|^2\to0$。\textbf{纠正：}课件的 $\varphi_n(z^q)$ 应为 $\varphi_n(x^q)$。

\Slide{15}{可行集、最优值、最优解是三个对象}
给定参数 $x$，可行行动是 $y\in\varphi(x)$，收益是连续函数 $f(x,y)$。定义
\[g(x)=\max_{y\in\varphi(x)}f(x,y),\qquad
\mu(x)=\{y\in\varphi(x):f(x,y)=g(x)\}.\]
$g(x)$ 是一个数，$\mu(x)$ 是一个集合。为保证它们有意义，必须假设 $\varphi(x)$ 非空且紧；本节共同紧 $T$ 下，闭值即可得紧值。极值定理说明最大值确实取到，故 $\mu(x)\ne\varnothing$。只有上确界存在不够，例如在 $(0,1)$ 最大化 $y$ 没有最优点。

\Slide{16}{图3：最优解集嵌在可行集里}
图中左侧是环境空间 $S$ 中的点 $x$，右侧依次是行动空间 $T$、可行集 $\varphi(x)$ 和最优解集 $\mu(x)$，表达
\[\mu(x)\subseteq\varphi(x)\subseteq T.\]
图中画成椭圆只是示意，不代表这些集合必为凸集、连通集或有内部。多个分开的点也可以同时最优；例如在 $[-1,1]$ 最大化 $y^2$，最优集是 $\{-1,1\}$。

\Slide{17}{Berge最大值定理的结论}
若收益 $f$ 连续、可行对应非空紧值且连续，则最优值 $g$ 连续、最优解对应 $\mu$ 非空紧值且上半连续。结论\emph{没有保证} $\mu$ 下半连续。课件括号中的“不要尝试证明 $g$ 连续”是原文的阅读安排，不是本自学讲义的限制；为适应证明考试，19页提供完整证明。

证明路线先记为三步：下半连续把旧的最优行动搬到附近，限制最优值不能骤降；上半连续和紧性把附近的最优行动抽取极限，限制最优值不能骤升；最后将最优等式取极限，保留最优性。

\Slide{18}{把最大值定理的假设写完整}
本页强调紧值性。准确的局部版本是：在 $x_0$ 附近每个 $\varphi(x)$ 非空紧，$\varphi$ 在 $x_0$ 上、下半连续（开集定义），$f$ 在 $S\times T$ 连续。结论为 $g$ 在 $x_0$ 连续，$\mu$ 在 $x_0$ 上半连续且各值非空紧。在本讲前段还可使用共同紧 $T$ 简化抽子列。

紧值只说“每个集合本身紧”；u.s.c. 负责“随着输入接近，集合不会逃走”。两者不能互相代替。课件 $\mu:S\to T$ 更严谨写成 $\mu:S\rightrightarrows T$，因为最优行动可以不唯一。

\Slide{19}{最优值连续的完整证明与原页纠错}
固定任意 $x_n\to x_0$，记 $g_0=g(x_0)$。\textbf{先证明下界。}取一个最优点 $y_0\in\mu(x_0)$。由 $\varphi$ 下半连续，可选 $z_n\in\varphi(x_n)$ 使 $z_n\to y_0$。连续性给
\[f(x_n,z_n)\to f(x_0,y_0)=g_0,\qquad g(x_n)\ge f(x_n,z_n),\]
故 $\liminf_n g(x_n)\ge g_0$。这里用 $\liminf$，因为还没有证明 $g(x_n)$ 的极限存在。

\textbf{再证明上界。}若 $\limsup_n g(x_n)>g_0$，则存在 $\eps>0$ 和子列，使 $g(x_{n_k})\ge g_0+\eps$。取\emph{最优点} $y_k\in\mu(x_{n_k})$，所以 $f(x_{n_k},y_k)=g(x_{n_k})$。共同紧 $T$ 可抽子列 $y_{k_j}\to y$，闭图给 $y\in\varphi(x_0)$。取极限得
\[f(x_0,y)=\lim_j f(x_{n_{k_j}},y_{k_j})\ge g_0+\eps,\]
但可行性要求 $f(x_0,y)\le g_0$，矛盾。于是 $\limsup_n g(x_n)\le g_0$，上下夹住得 $g(x_n)\to g_0$。一般紧值 u.s.c. 情形用第87页的子列结论代替共同紧 $T$。

\textbf{纠正：}原页把未知是否存在的 $\lim g(x_n)$ 直接写出来，并把任意可行 $y_n$ 当成最优点；还过早写 $g(x_0)=f(x_0,y_0)$。以上证明逐项补上所需理由。

\Slide{20}{最优解对应为何紧且上半连续}
对固定 $x$，极值定理给最优点存在。$\mu(x)=\varphi(x)\cap\{y:f(x,y)=g(x)\}$ 是紧集的闭子集，所以紧。接着任取 $x_n\to x_0$、$y_n\in\mu(x_n)$、$y_n\to y$。因为 $\mu(x_n)\subseteq\varphi(x_n)$，闭图给 $y\in\varphi(x_0)$；又
\[f(x_n,y_n)=g(x_n)\ \Longrightarrow\ f(x_0,y)=g(x_0),\]
故 $y\in\mu(x_0)$，最优图闭。

在共同紧 $T$ 中闭图即上半连续。一般情形：若某开集 $G\supseteq\mu(x_0)$ 不能包含附近全部最优点，就可选 $x_n\to x_0$、$y_n\in\mu(x_n)\setminus G$；紧值 u.s.c. 的 $\varphi$ 给收敛子列，刚才论证使其极限在 $\mu(x_0)\subseteq G$，但各点在闭集 $T\setminus G$ 中，矛盾。只证明“某个好子列存在”而不说明怎样排除坏序列，容易留下逻辑缺口。

\Slide{21}{最优值连续而最优解不下半连续}
$S=T=[0,1]$，$\varphi(x)=[0,1]$，$f(x,y)=xy$。$x>0$ 时乘以正数保序，所以唯一最优点为1；$x=0$ 时所有收益都为0，故 $\mu(0)=[0,1]$。$g(x)=x$ 连续。

$\mu$ 在0上半连续，因为任何附近最优点的极限都在 $[0,1]$。但不下半连续：取 $x_n=1/n$ 和目标 $y_0=0\in\mu(0)$，每个 $y_n\in\mu(x_n)$ 都只能等于1，不能趋于0。这个反例说明最大值定理的结论已经是恰当强度，不能擅自把 u.s.c. 改为连续。

\Slide{22}{紧值域假设什么时候能替换}
若行动空间原本是 $\R^L$，可以证明所有相关可行点（或所有相关最优点）位于某个共同紧集，再在该集上运用前面的定理。例如收入 $w$ 局部有上界、每个商品价格 $p_i$ 局部有共同正下界 $c>0$，预算约束推出 $0\le y_i\le w/c$，从而有统一有界容器。价格允许趋于零时这个论证会失效。不可直接把不紧的 $T$ 写成紧而不证明截断不会丢掉所需选择。

\Slide{23}{上下半连续对应的邻域含义}
下半连续：每个与 $\Phi(x)$ 相交的开集 $V$，在附近每个输入的值集仍与 $V$ 相交。上半连续：每个包含整个 $\Phi(x)$ 的开集 $V$，附近全部值集都包含于 $V$。前者保证旧选择可近似保留，后者禁止新选择逃出旧值集的开邻域。对单点值，两定义均化为普通函数连续；这里的上、下不是实值函数大小方向。

\Slide{24}{上下半连续对应的邻域含义}
本页继续“上下半连续对应的邻域含义”。完整定义、条件和论证见本讲原页第 23 页下方。

\Slide{25}{上下半连续对应的邻域含义}
本页继续“上下半连续对应的邻域含义”。完整定义、条件和论证见本讲原页第 23 页下方。

\Slide{26}{上下半连续对应的邻域含义}
本页继续“上下半连续对应的邻域含义”。完整定义、条件和论证见本讲原页第 23 页下方。

\Slide{27}{上下半连续对应的邻域含义}
本页继续“上下半连续对应的邻域含义”。完整定义、条件和论证见本讲原页第 23 页下方。

\Slide{28}{逆像刻画的证明}
下逆像 $\Phi^-(V)=\{x:\Phi(x)\cap V\ne\varnothing\}$。下半连续当且仅开集的下逆像开：若某点属于该逆像，定义给整个邻域也属于它；反向取逆像本身为邻域。上逆像 $\Phi^+(V)=\{x:\Phi(x)\subseteq V\}$，同样论证得到上半连续的开逆像判据。非空值下的补集对偶把两种开集判据换成闭集判据；允许空值时应先限制到有效定义域。

\Slide{29}{逆像刻画的证明}
本页继续“逆像刻画的证明”。完整定义、条件和论证见本讲原页第 28 页下方。

\Slide{30}{逆像刻画的证明}
本页继续“逆像刻画的证明”。完整定义、条件和论证见本讲原页第 28 页下方。

\Slide{31}{逆像刻画的证明}
本页继续“逆像刻画的证明”。完整定义、条件和论证见本讲原页第 28 页下方。

\Slide{32}{逆像刻画的证明}
本页继续“逆像刻画的证明”。完整定义、条件和论证见本讲原页第 28 页下方。

\Slide{33}{逆像刻画的证明}
本页继续“逆像刻画的证明”。完整定义、条件和论证见本讲原页第 28 页下方。

\Slide{34}{紧值上半连续对应的紧像}
设 $K$ 紧且 $\Phi$ 紧值上半连续。给 $\Phi(K)$ 的任意开覆盖，对每个 $x\in K$ 先有限覆盖紧集 $\Phi(x)$，把对应有限个开集的并记作 $V_x$。上半连续给 $U_x$ 使 $\Phi(U_x)\subset V_x$。再用 $K$ 紧，从 $U_x$ 取有限子覆盖。两次有限选择的总开集数仍有限，覆盖 $\Phi(K)$。两个紧性分别作用于“每个值”和“输入集”。

\Slide{35}{闭对应与闭图}
闭图指若 $x_n\to x,y_n\to y,y_n\in\Phi(x_n)$，则 $y\in\Phi(x)$。等价于图在乘积空间闭；也等价于图外每个 $(x,y)$ 可找到乘积开邻域与图不交。固定 $x$ 用常数输入得闭图蕴含闭值。连续 $f(x,y)$ 的弱不等式定义对应有闭图，因为不等式保极限。任意闭图对应的逐点交为图的任意闭交，有限并为有限闭并，有限积由坐标判据闭。闭值本身不推出闭图，例如单值跳跃函数。

\Slide{36}{闭对应与闭图}
本页继续“闭对应与闭图”。完整定义、条件和论证见本讲原页第 35 页下方。

\Slide{37}{闭对应与闭图}
本页继续“闭对应与闭图”。完整定义、条件和论证见本讲原页第 35 页下方。

\Slide{38}{闭对应与闭图}
本页继续“闭对应与闭图”。完整定义、条件和论证见本讲原页第 35 页下方。

\Slide{39}{闭对应与闭图}
本页继续“闭对应与闭图”。完整定义、条件和论证见本讲原页第 35 页下方。

\Slide{40}{闭对应与闭图}
本页继续“闭对应与闭图”。完整定义、条件和论证见本讲原页第 35 页下方。

\Slide{41}{闭对应与闭图}
本页继续“闭对应与闭图”。完整定义、条件和论证见本讲原页第 35 页下方。

\Slide{42}{闭对应与闭图}
本页继续“闭对应与闭图”。完整定义、条件和论证见本讲原页第 35 页下方。

\Slide{43}{闭对应与闭图}
本页继续“闭对应与闭图”。完整定义、条件和论证见本讲原页第 35 页下方。

\Slide{44}{闭对应与闭图}
本页继续“闭对应与闭图”。完整定义、条件和论证见本讲原页第 35 页下方。

\Slide{45}{闭图与上半连续之间的紧性桥梁}
紧值上半连续对应有闭图：若 $y\notin\Phi(x)$，紧值闭且 $\dist(y,\Phi(x))>0$，取分离开邻域，再用上半连续，附近图不能靠近 $(x,y)$。若 $\Psi$ 闭图，$\Phi$ 紧值上半连续，则 $\Phi\cap\Psi$ 紧值上半连续。证明反证：若有 $x_n\to x$ 和交值 $y_n$ 逃出包含交值的开集 $G$，紧像性质给子列 $y_n\to y$，两图闭使 $y$ 属于交值，开性使子列最终进 $G$，矛盾。值域空间紧时，闭图亦推出上半连续：逃出开集的选择总能抽取收敛子列。

\Slide{46}{闭图与上半连续之间的紧性桥梁}
本页继续“闭图与上半连续之间的紧性桥梁”。完整定义、条件和论证见本讲原页第 45 页下方。

\Slide{47}{闭图与上半连续之间的紧性桥梁}
本页继续“闭图与上半连续之间的紧性桥梁”。完整定义、条件和论证见本讲原页第 45 页下方。

\Slide{48}{闭图与上半连续之间的紧性桥梁}
本页继续“闭图与上半连续之间的紧性桥梁”。完整定义、条件和论证见本讲原页第 45 页下方。

\Slide{49}{闭图与上半连续之间的紧性桥梁}
本页继续“闭图与上半连续之间的紧性桥梁”。完整定义、条件和论证见本讲原页第 45 页下方。

\Slide{50}{闭图与上半连续之间的紧性桥梁}
本页继续“闭图与上半连续之间的紧性桥梁”。完整定义、条件和论证见本讲原页第 45 页下方。

\Slide{51}{闭图与上半连续之间的紧性桥梁}
本页继续“闭图与上半连续之间的紧性桥梁”。完整定义、条件和论证见本讲原页第 45 页下方。

\Slide{52}{紧空间上的不变紧集}
假设非空紧 $S$、非空紧值上半连续自对应。令 $K_n=\Phi^n(S)$，紧像定理使它们非空紧，单调性使 $K_{n+1}\subseteq K_n$，故交 $K$ 非空紧。单调性给 $\Phi(K)\subseteq K$。对 $a\in K$，可取 $x_n\in K_n$ 使 $a\in\Phi(x_n)$，抽子列 $x_n\to x$，套集闭性给 $x\in K$，闭图给 $a\in\Phi(x)$。故 $K=\Phi(K)$。这不保证某点 $x\in\Phi(x)$，不要把不变集合误当不动点。

\Slide{53}{紧空间上的不变紧集}
本页继续“紧空间上的不变紧集”。完整定义、条件和论证见本讲原页第 52 页下方。

\Slide{54}{紧空间上的不变紧集}
本页继续“紧空间上的不变紧集”。完整定义、条件和论证见本讲原页第 52 页下方。

\Slide{55}{复合与下半连续并}
复合的下逆像满足 $(\Psi\circ\Phi)^-(V)=\Phi^-(\Psi^-(V))$，故两次下半连续使开集逐次变开；上逆像的全称量词也能逐层写出，得上半连续复合。任意并的下逆像为各下逆像之并，因此任意并保下半连续。每个等式都来自“存在中间选择”或“所有中间选择”，而不是把对应当单值函数计算。

\Slide{56}{复合与下半连续并}
本页继续“复合与下半连续并”。完整定义、条件和论证见本讲原页第 55 页下方。

\Slide{57}{复合与下半连续并}
本页继续“复合与下半连续并”。完整定义、条件和论证见本讲原页第 55 页下方。

\Slide{58}{复合与下半连续并}
本页继续“复合与下半连续并”。完整定义、条件和论证见本讲原页第 55 页下方。

\Slide{59}{复合与下半连续并}
本页继续“复合与下半连续并”。完整定义、条件和论证见本讲原页第 55 页下方。

\Slide{60}{复合与下半连续并}
本页继续“复合与下半连续并”。完整定义、条件和论证见本讲原页第 55 页下方。

\Slide{61}{对应交集的条件与伪定理}
紧值上半连续对应的任意交是闭图，固定一个紧值上半连续因子作容器，运用第45页交集定理得紧值上半连续（交值可能为空）。第63--64页删掉紧值后的证明不能沿用，尤其 $(\bigcap_i\Phi_i)^-(F)$ 一般不等于 $\bigcap_i\Phi_i^-(F)$：各因子与 $F$ 相交可能在不同点。第65--66页明确为伪定理：下半连续的有限交不必下半连续。反例 $\Phi_1(x)=\{0\}$，$\Phi_2(x)=\{x\}$，交仅在0非空；若要求非空值，可改 $\Phi_1(x)=\{0,1\},\Phi_2(x)=\{x,1\}$，在0交含0与1，非零附近只含1，不下半连续。

\Slide{62}{对应交集的条件与伪定理}
本页继续“对应交集的条件与伪定理”。完整定义、条件和论证见本讲原页第 61 页下方。

\Slide{63}{对应交集的条件与伪定理}
本页继续“对应交集的条件与伪定理”。完整定义、条件和论证见本讲原页第 61 页下方。

\Slide{64}{对应交集的条件与伪定理}
本页继续“对应交集的条件与伪定理”。完整定义、条件和论证见本讲原页第 61 页下方。

\Slide{65}{对应交集的条件与伪定理}
本页继续“对应交集的条件与伪定理”。完整定义、条件和论证见本讲原页第 61 页下方。

\Slide{66}{对应交集的条件与伪定理}
本页继续“对应交集的条件与伪定理”。完整定义、条件和论证见本讲原页第 61 页下方。

\Slide{67}{有限并保持上半连续}
若开集 $G$ 包含 $\bigcup_{i=1}^m\Phi_i(x)$，它包含每个 $\Phi_i(x)$。为每个因子取邻域 $U_i$ 使其值在 $G$ 内，有限交 $U$ 仍为邻域，故并值在 $G$ 内。紧值时有限并亦紧。原页逆像公式须注意：并对应的上逆像是各上逆像的交，下逆像是各下逆像的并。不能把两者写成同一种交并。

\Slide{68}{有限并保持上半连续}
本页继续“有限并保持上半连续”。完整定义、条件和论证见本讲原页第 67 页下方。

\Slide{69}{有限并保持上半连续}
本页继续“有限并保持上半连续”。完整定义、条件和论证见本讲原页第 67 页下方。

\Slide{70}{有限并保持上半连续}
本页继续“有限并保持上半连续”。完整定义、条件和论证见本讲原页第 67 页下方。

\Slide{71}{有限积对应的半连续性}
下半连续积：给定值点 $y=(y_i)$ 和含它的开集，取乘积开邻域 $\prod_iV_i$ 包含于其中。各因子的下半连续保证附近可分别选 $y_i'\in\Phi_i(x')\cap V_i$，于是组成积中选择。紧值上半连续积：给任意输入收敛序列及积值选择，依次抽每个因子的收敛子列，有限次后所有坐标同时收敛且极限在积值中；若不满足上半连续，选逃离开邻域的序列即矛盾。没有紧值的第73--74页不能用上述紧性证明，也不能把任意开集当单个开矩形；本讲使用有紧值条件的正确版本。

\Slide{72}{有限积对应的半连续性}
本页继续“有限积对应的半连续性”。完整定义、条件和论证见本讲原页第 71 页下方。

\Slide{73}{有限积对应的半连续性}
本页继续“有限积对应的半连续性”。完整定义、条件和论证见本讲原页第 71 页下方。

\Slide{74}{有限积对应的半连续性}
本页继续“有限积对应的半连续性”。完整定义、条件和论证见本讲原页第 71 页下方。

\Slide{75}{有限积对应的半连续性}
本页继续“有限积对应的半连续性”。完整定义、条件和论证见本讲原页第 71 页下方。

\Slide{76}{有限积对应的半连续性}
本页继续“有限积对应的半连续性”。完整定义、条件和论证见本讲原页第 71 页下方。

\Slide{77}{和对应与凸包}
有限和是积对应经连续加法映射的像，故紧值上半连续性由积与复合得到。下半连续凸包：任取值中一点 $y=\sum_{i=1}^m\lambda_i y_i$，对每个 $y_i$ 用下半连续选取随输入趋近的 $y_{in}$，保持同一组权重，便得凸组合趋于 $y$。紧值上半连续凸包在 $\R^N$ 也成立：Carathéodory 定理把每点表示为至多 $N+1$ 个点的凸组合，权重在紧单纯形中；对权重和值点抽子列得极限仍在凸包。有限维假设不可遗漏。

\Slide{78}{和对应与凸包}
本页继续“和对应与凸包”。完整定义、条件和论证见本讲原页第 77 页下方。

\Slide{79}{和对应与凸包}
本页继续“和对应与凸包”。完整定义、条件和论证见本讲原页第 77 页下方。

\Slide{80}{和对应与凸包}
本页继续“和对应与凸包”。完整定义、条件和论证见本讲原页第 77 页下方。

\Slide{81}{和对应与凸包}
本页继续“和对应与凸包”。完整定义、条件和论证见本讲原页第 77 页下方。

\Slide{82}{闭包对应与内半连续}
第82--83页的横线运算是闭包，不应标为 convex hull。$\Phi$ 与逐点闭包的下逆像对开集相同，因为开集遇到闭包就遇到原集合，故下半连续等价。上半连续的闭包版本用正则性：给开集 $G$ 包含闭包值，先取开邻域 $V$ 使 $\Phi(x)\subset V\subset\overline V\subset G$（紧值时可有限选取），再用上半连续。度量空间中下半连续等价于：任意 $x_n\to x,y\in\Phi(x)$ 可选 $y_n\in\Phi(x_n)$ 趋于 $y$。证明依次对球 $B(y,1/k)$ 取输入邻域，再用足够大的指标分段选择。

\Slide{83}{闭包对应与内半连续}
本页继续“闭包对应与内半连续”。完整定义、条件和论证见本讲原页第 82 页下方。

\Slide{84}{闭包对应与内半连续}
本页继续“闭包对应与内半连续”。完整定义、条件和论证见本讲原页第 82 页下方。

\Slide{85}{闭包对应与内半连续}
本页继续“闭包对应与内半连续”。完整定义、条件和论证见本讲原页第 82 页下方。

\Slide{86}{闭包对应与内半连续}
本页继续“闭包对应与内半连续”。完整定义、条件和论证见本讲原页第 82 页下方。

\Slide{87}{外半连续的术语与序列选取}
本课此处的 o.s.c. 要求：任意 $x_n\to x$ 及 $y_n\in\Phi(x_n)$，有收敛子列且极限属于 $\Phi(x)$。它比文献中仅要求“已有收敛选择的极限不逃出”的闭图式 outer semicontinuity 更强。紧值上半连续给该性质：集合 $\{x,x_1,x_2,\ldots\}$ 紧，紧像定理给子列，闭图给极限归属。反向若上半连续失效，选开 $G\supset\Phi(x)$ 及逃出 $G$ 的 $y_n$，由假设抽子列极限在 $\Phi(x)$，与开性矛盾。复合用先抽中间值再抽末端值，有限积用有限次子列抽取；不能无限次抽取后还把未构造的子列当作存在。

\Slide{88}{外半连续的术语与序列选取}
本页继续“外半连续的术语与序列选取”。完整定义、条件和论证见本讲原页第 87 页下方。

\Slide{89}{外半连续的术语与序列选取}
本页继续“外半连续的术语与序列选取”。完整定义、条件和论证见本讲原页第 87 页下方。

\Slide{90}{外半连续的术语与序列选取}
本页继续“外半连续的术语与序列选取”。完整定义、条件和论证见本讲原页第 87 页下方。

\Slide{91}{外半连续的术语与序列选取}
本页继续“外半连续的术语与序列选取”。完整定义、条件和论证见本讲原页第 87 页下方。

\Slide{92}{外半连续的术语与序列选取}
本页继续“外半连续的术语与序列选取”。完整定义、条件和论证见本讲原页第 87 页下方。

\Slide{93}{外半连续的术语与序列选取}
本页继续“外半连续的术语与序列选取”。完整定义、条件和论证见本讲原页第 87 页下方。

\Slide{94}{外半连续的术语与序列选取}
本页继续“外半连续的术语与序列选取”。完整定义、条件和论证见本讲原页第 87 页下方。

\Slide{95}{外半连续的术语与序列选取}
本页继续“外半连续的术语与序列选取”。完整定义、条件和论证见本讲原页第 87 页下方。

\Slide{96}{最优值的两种半连续性}
设 $M(x)=\sup_{y\in\Phi(x)}f(x,y)$。若 $f$ 下半连续且 $\Phi$ 下半连续，任取 $a<M(x)$，选可行 $y$ 满足 $f(x,y)>a$。目标下半连续给乘积邻域使函数值仍大于 $a$，可行对应下半连续给附近可行点落入该邻域，故附近 $M>a$。若 $f$ 上半连续，$\Phi$ 非空紧值上半连续，则每点最大值存在。取输入序列及最优点，由紧像抽子列趋于可行极限，上半连续给 $\limsup M(x_n)\le M(x)$。对 limsup 严格写法需先选逼近它的子列，再抽最优点子列。

\Slide{97}{最优值的两种半连续性}
本页继续“最优值的两种半连续性”。完整定义、条件和论证见本讲原页第 96 页下方。

\Slide{98}{最优值的两种半连续性}
本页继续“最优值的两种半连续性”。完整定义、条件和论证见本讲原页第 96 页下方。

\Slide{99}{Berge 最大值与唯一选择}
必须补充非空紧值：$\Phi$ 连续且非空紧值，$f$ 连续，则第96页的两个方向给 $M$ 连续。最优对应 $A(x)=\{y\in\Phi(x):f(x,y)=M(x)\}$ 非空紧；等式由连续性保极限，故闭图，再以 $\Phi$ 为紧容器得上半连续。若最优点唯一，单点值上半连续等价普通连续，故唯一最优选择连续。实数紧值取 $f(y)=y$，最大值唯一，得到连续选择；没有紧值时例如 $(0,1)$ 根本没有最大点。原页100的球应理解为非空闭球，开球一般不取到坐标最大值。

\Slide{100}{Berge 最大值与唯一选择}
本页继续“Berge 最大值与唯一选择”。完整定义、条件和论证见本讲原页第 99 页下方。

\Slide{101}{Berge 最大值与唯一选择}
本页继续“Berge 最大值与唯一选择”。完整定义、条件和论证见本讲原页第 99 页下方。

\Slide{102}{Berge 最大值与唯一选择}
本页继续“Berge 最大值与唯一选择”。完整定义、条件和论证见本讲原页第 99 页下方。

\Slide{103}{最小化的对偶证明}
对最小化使用 $-f$ 的最大化，或直接证明：紧值上半连续可行集及下半连续目标给最小值函数下半连续；下半连续可行集及上半连续目标给下确界函数上半连续。后一证明先选 $y_0$ 使 $f(x_0,y_0)<m(x_0)+\eps/2$，再使附近某可行选择的目标小于 $f(x_0,y_0)+\eps/2$。两项合成 $m(x)<m(x_0)+\eps$。第104页末尾应是 $M(x)$，不是 $M(x_0)$；第106页错误的集合包含应改成存在 $y\in\Phi(x)\cap V$。第107页最小对应的等式最安全写为 $f(y)=m(x)$，原图弱不等式方向不能把所有较大值也选入。

\Slide{104}{最小化的对偶证明}
本页继续“最小化的对偶证明”。完整定义、条件和论证见本讲原页第 103 页下方。

\Slide{105}{最小化的对偶证明}
本页继续“最小化的对偶证明”。完整定义、条件和论证见本讲原页第 103 页下方。

\Slide{106}{最小化的对偶证明}
本页继续“最小化的对偶证明”。完整定义、条件和论证见本讲原页第 103 页下方。

\Slide{107}{最小化的对偶证明}
本页继续“最小化的对偶证明”。完整定义、条件和论证见本讲原页第 103 页下方。

\Slide{108}{环境依赖偏好与最大元}
假设 $\beta$ 非空紧值连续，每个严格关系传递且非自反，且“可行且非严格优于”的三元组集合闭。若无极大元，每个 $y$ 被某个 $z$ 严格压过，这些严格下截集构成 $\beta(x)$ 的相对开覆盖。有限子覆盖的有限个 $z_i$ 中存在极大者（传递且无循环），它不被任一 $z_i$ 压过，矛盾。极大元集为所有非优于截集的交，故闭且紧。若 $x_n\to x,y_n$ 为极大元且 $y_n\to y$，先由可行图闭得 $y\in\beta(x)$，再对任意 $z\in\beta(x)$ 用下半连续选 $z_n\in\beta(x_n)\to z$。由 $y_n$ 不差于 $z_n$ 及三元组闭性得极限不被 $z$ 压过。第111页 $z_n$ 应属于可行集 $\beta(x_n)$，不必属于最优集。以连续效用定义严格偏好即可回到 Berge 定理；第112页最优对应内目标须写成 $f(x,y)$，并补紧值条件。

\Slide{109}{环境依赖偏好与最大元}
本页继续“环境依赖偏好与最大元”。完整定义、条件和论证见本讲原页第 108 页下方。

\Slide{110}{环境依赖偏好与最大元}
本页继续“环境依赖偏好与最大元”。完整定义、条件和论证见本讲原页第 108 页下方。

\Slide{111}{环境依赖偏好与最大元}
本页继续“环境依赖偏好与最大元”。完整定义、条件和论证见本讲原页第 108 页下方。

\Slide{112}{环境依赖偏好与最大元}
本页继续“环境依赖偏好与最大元”。完整定义、条件和论证见本讲原页第 108 页下方。

\Slide{113}{环境依赖偏好与最大元}
本页继续“环境依赖偏好与最大元”。完整定义、条件和论证见本讲原页第 108 页下方。

\clearpage\section*{补充：凸包所用的 Carathéodory 定理}
在 $\R^N$，每个凸包点都能用至多 $N+1$ 个原集合点表示。证明：从任意有限表示 $x=\sum_{i=1}^m\lambda_i x_i$ 开始，去掉零权重。若 $m>N+1$，向量 $(x_i,1)\in\R^{N+1}$ 线性相关，故有不全零 $a_i$ 满足 $\sum_i a_ix_i=0,\sum_i a_i=0$。正负 $a_i$ 均存在，取 $t=\min_{a_i>0}\lambda_i/a_i>0$，新权重 $\lambda_i-ta_i$ 均非负、和仍1、表示的点不变，且至少一个权重变0。去掉该点，有限次后至多剩 $N+1$ 点。因此紧集凸包是紧单纯形乘以 $K^{N+1}$ 在连续凸组合映射下的像，仍紧。
\paragraph{证明自测} 为什么最大值定理只给最优对应上半连续？例目标 $f(x,y)=xy$，可行集恒为 $[-1,1]$。$x=0$ 时全部点最优，$x>0$ 时仅1最优，$x<0$ 时仅 $-1$ 最优；原点的旧最优点0无法由正参数最优点逼近，所以下半连续失败。最优值为 $|x|$，仍连续。
